\documentclass[11pt]{article}
\usepackage[letterpaper,textwidth=5.8in,textheight=8.8in]{geometry}
\usepackage{amsmath,amssymb,graphicx,booktabs,array,tabularx}
\usepackage[numbers,sort&compress]{natbib}
\usepackage[hidelinks]{hyperref}
\usepackage{microtype}
\usepackage[font=small,labelfont=bf]{caption}
\setlength{\emergencystretch}{2em}
\newcommand{\figpath}{../figures}
\title{Prediction Is Not Yet Control:\\An Empirical Audit of a Small World--Action Model}
\author{Keivalya Pandya\\Northeastern University}
\date{}
\begin{document}
\maketitle
\begin{abstract}
A learned forward model can predict demonstrations accurately without supporting reliable decisions about alternative actions. We examine this distinction in mini-WAM, a compact policy with shared latent dynamics, using a two-dimensional navigation environment whose transitions can be checked directly. Archived experiments across three independent runs show low state-prediction error, but no policy improvement when the physical state is supplied. Removing state inputs increases the observed success gap in favor of joint prediction, from zero to 44.3 percentage points, although privileged state supervision and unequal training windows prevent a clean attribution to dynamics learning. A fresh five-run supervision pilot obtains 17.0\% success with behavioral cloning, 28.4\% with current-state supervision, and 33.6\% with additional future-state prediction. Image predictions also perform worse than copying the current frame. A fresh five-run system-identification audit finds that an action-conditioned affine model fits the toy dynamics to numerical precision, even with deterministic expert data. This result narrows the interpretation of the neural model's counterfactual failures: they are not an unavoidable property of this environment or its demonstrations. We separate predictive accuracy, observation sufficiency, and control utility, and identify the comparisons needed before claiming that learned world models improve decision-making.
\end{abstract}

\section{Introduction}
A policy and a forward model answer different questions. A policy selects an action from the information available now. A forward model predicts the consequences of an action, including actions the policy would rarely choose. Sharing their representations is appealing: prediction may encourage a policy to encode features that matter for control. Yet a reduction in prediction error does not establish that the resulting policy makes better decisions.

The distinction is easy to lose in a small demonstration. A model may produce plausible rollouts, expose interpretable latent coordinates, and reach a goal reliably. Each observation is useful, but none by itself shows that the forward model is necessary. If the state already specifies the goal and velocity, a reactive controller may solve the task. If evaluation follows the expert's actions, a predictor may exploit regularities of the expert rather than learn a transition rule that generalizes to other controls.

We study these issues in mini-WAM, an existing compact implementation of a policy and action-conditioned latent predictor. The environment is deliberately simple: a point mass moves with damped momentum toward one of two colored goals. Its simplicity permits exact interventions and inexpensive comparisons. It also makes strong baselines essential. Away from position clipping at the arena boundary, its dynamics are affine; a large improvement over a stationary prediction is therefore a weak test of modeling ability.

This paper asks three questions. First, how does prediction on recorded trajectories relate to prediction under alternative actions? Second, what can joint prediction explain about policy performance when observations omit velocity? Third, which apparent benefits survive comparison with a model matched to the environment's structure?

The evidence gives a qualified answer. The archived neural results contain a substantial gap between expert-only and exploratory data in counterfactual prediction. They also contain a promising, variable improvement in pixel-input control. However, a new affine baseline learns the transition law from expert trajectories alone, and a direct rendering test confirms that opposite velocities can produce identical observations. These findings rule out two stronger interpretations: that exploration is intrinsically necessary to identify this particular simulator, and that imagined rollouts recover information missing from a single image without further assumptions.

Our contribution is an empirical audit, rather than a new world-model architecture. We bring the archived measurements, executable simulator checks, and missing controls into one account. The scope is a single toy task family. We do not claim a new planning algorithm, general language grounding, or evidence of transfer to physical robots.

\section{Related work}
Learning a compact representation and predicting its evolution is established practice in model-based reinforcement learning. World Models \citep{ha2018} and Dreamer \citep{hafner2020} use learned dynamics to support behavior learning. TD-MPC2 \citep{hansen2024} couples latent prediction to online planning. mini-WAM shares their use of latent transitions, but its evaluated policy does not plan or optimize behavior inside the model. Its forward model supplies an auxiliary training signal and an object of analysis.

A second line of work combines action prediction with representations of future observations. Unified World Models \citep{zhu2025} couples video and action diffusion. The recent, unrelated MiniWAM \citep{chen2026} learns compact predictive targets for robot policies. Our repository's name predates this comparison, but the systems should not be conflated: mini-WAM here is a small supervised diagnostic implementation, not that method or a reproduction of its results. Compactness alone is not an architectural novelty claim.

The concern about factual prediction is also not new. Causal confusion in imitation learning \citep{dehaan2019} shows how observational regularities can mislead policies. More directly, AD-WM \citep{qiu2026} studies action discrimination for counterfactual model predictive control. It proposes action-recovery objectives and evaluates decision-relevant diagnostics. Consequently, observing a factual--counterfactual gap in mini-WAM does not establish priority for that problem. A stronger contribution would need a distinct mechanism, a controlled evaluation protocol with demonstrated value, or evidence that existing explanations fail in a meaningful setting.

Finally, our linear state head is supervised jointly with the representation. Its accuracy measures whether the requested state can be read out after training; it is not an independent probe of spontaneously learned structure. Post-hoc probes \citep{alain2016} require a separate fitting and evaluation procedure. This distinction matters especially when the same state variables are both encoder inputs and reconstruction targets.

\section{Environment and observation model}
\label{sec:environment}
The agent moves in the unit square. Each episode contains a blue goal and a green goal, and an instruction selects one of them. The renderer produces a $64\times64$ RGB image with a red disk for the agent and colored squares for the goals (Figure~\ref{fig:env}). Initial velocity is zero. Agent and goal positions are sampled with minimum separations so that episodes do not begin at success.

\begin{figure}[t]
\centering
\includegraphics[width=\linewidth]{\figpath/research-environment.pdf}
\caption{Actual simulator frames from a deterministic expert rollout. The instruction selects the blue square; the green square remains a distractor. Positions and colors are rendered by the environment, without illustrative replacements. Velocity is not drawn. These frames describe the task, not a learned policy's performance.}
\label{fig:env}
\end{figure}

Let $p_t\in[0,1]^2$ be position, $v_t\in\mathbb R^2$ physical velocity, and $a_t\in[-1,1]^2$ the applied force. The transition is
\begin{equation}
 v_{t+1}=0.8v_t+0.01a_t,\qquad
 p_{t+1}=\operatorname{clip}(p_t+v_{t+1},0,1).
 \label{eq:physics}
\end{equation}
The state returned by the environment is $s_t=(p_t,u_t,g)$, where $u_t=v_t/0.05$ is normalized velocity and $g$ is the selected goal position. Thus $u_{t+1}=0.8u_t+0.2a_t$. Success requires distance to the selected goal below 0.05 and normalized speed below 0.2. Position clipping does not reset velocity; this boundary rule must be preserved by counterfactual simulations.

The expert is a clipped proportional--derivative controller,
\begin{equation}
 a_t=\operatorname{clip}\left(5(g-p_t)-2u_t+\epsilon_t,-1,1\right),
 \qquad \epsilon_t\sim\mathcal N(0,0.1^2 I).
\end{equation}
The collection procedure can replace expert actions with short random-action bursts. Those transitions train the world model but are masked out of the imitation loss. Burst-start probability is a collection parameter, not the measured fraction of random transitions: burst duration and episode termination also affect that fraction.

Two observation regimes must be distinguished. In the state-input regime, the encoder receives the image, instruction, and all six state coordinates. Because $g$ already identifies the selected goal, this regime does not require language understanding or visual localization. In the pixel-input regime, the state-encoder contribution is zeroed out. The WAM nevertheless retains supervised state-prediction targets during training. It is therefore a pixel-input policy with privileged training supervision, not a method trained from pixels alone.

The renderer depends on positions and goal colors, but not velocity. Holding the scene fixed while reversing velocity produces exactly the same image. A memoryless encoder therefore cannot distinguish those states. Correlations in the training distribution may still make velocity partly predictable, but they do not resolve the underlying observation ambiguity.

\section{Model and learning objective}
The encoder combines a spatial-softmax image representation, a small text GRU, and a state MLP into a 128-dimensional latent vector. A fusion MLP followed by layer normalization produces
\begin{equation}
 z_t=E_\theta(I_t,\ell,\tilde s_t),\qquad
 a_t=\pi_\theta(z_t),
\end{equation}
where $\tilde s_t$ uses training-set normalization. In the pixel-input condition the state branch is replaced by zeros before fusion. All reported archived experiments use the MLP action head; available diffusion and flow heads are not evaluated here.

The forward model applies a residual MLP to the latent and action:
\begin{equation}
 \hat z_{t+k+1}=F_\theta(\hat z_{t+k},a_{t+k}),\quad
 F_\theta(z,a)=\operatorname{LN}(z+f_\theta([z,a])),\quad \hat z_t=z_t.
\end{equation}
A linear head $R(z)=Wz+b$ predicts normalized state. An image decoder $D$ predicts RGB residuals around the training-set mean image. The same readouts operate on encoded and predicted latents. Layer normalization controls latent scale, but does not guarantee that predicted latents lie on the distribution of encoded observations.

For an $H$-step training window, define
\begin{align}
 \mathcal L_s&=\operatorname{MSE}(R(z_t),\tilde s_t)
 +\frac1H\sum_{k=1}^H\operatorname{MSE}(R(\hat z_{t+k}),\tilde s_{t+k}),\\
 \mathcal L_I&=\operatorname{MSE}(D(z_t),I_t)
 +\frac1H\sum_{k=1}^H\operatorname{MSE}(D(\hat z_{t+k}),I_{t+k}),\\
 \mathcal L_z&=\frac1H\sum_{k=1}^H\operatorname{MSE}\left(\hat z_{t+k},
 \operatorname{sg}[E_\theta(I_{t+k},\ell,\tilde s_{t+k})]\right).
\end{align}
The objective is $\mathcal L=\mathcal L_a+0.5\mathcal L_s+\mathcal L_I+\mathcal L_z$, with expert-only action regression $\mathcal L_a$ and stop-gradient operator $\operatorname{sg}$. Each MSE averages over its output coordinates. The default horizon is five. Removing all three auxiliary losses gives the behavioral-cloning comparator, called VLA in the repository. It has the same encoder design but is trained from scratch and should not be confused with a pretrained robot foundation model.

At deployment, both policies encode the current observation and emit an action. The WAM does not feed an imagined trajectory into its action head, retain a recurrent belief, or search over candidate plans. Any control difference in the archived experiments must arise through training, not online use of the forward model.

\section{Evidence and evaluation protocol}
\label{sec:protocol}
We distinguish two sources of evidence. The archived study consists of committed per-run JSON metrics and exported arrays used for figures. Its standard configuration uses 400 episodes, a 90/10 episode split, three independent training runs, 15 training epochs, and 100 closed-loop evaluation episodes per run. The initial audit checkout contains neither the original training datasets nor checkpoints. The archived numbers can be reaggregated, but that does not independently reproduce the training runs. Fresh simulator and affine-baseline experiments are described separately in Section~\ref{sec:audit}.

Offline neural evaluation rolls the latent forward from each valid start in a held-out episode, using the recorded actions without re-encoding future observations. At horizon $k$, state error is
\begin{equation}
 \operatorname{NMSE}_k=\frac{1}{N_kd}\sum_{i=1}^{N_k}\sum_{j=1}^{d}
 \left(\frac{\hat s^{(i)}_{t+k,j}-s^{(i)}_{t+k,j}}{\sigma_j}\right)^2,
\end{equation}
where $\sigma_j$ is the training-set standard deviation. Factual evaluation uses all six state dimensions; counterfactual evaluation uses only position and normalized velocity. Their values should not be compared as though they were the same metric.

For counterfactual evaluation, the simulator branches from held-out states under sampled random action sequences. The learned prediction is compared with that branch, and a zero-action simulator trajectory supplies a reference. This reference has access to the true initial state and dynamics; it tests the cost of disregarding the candidate action, not the best achievable learned predictor. The factual copy baseline instead holds the current state or image fixed.

Reported archived spreads are population standard deviations across three training runs, matching the saved aggregation code. They are descriptive spreads, not confidence intervals. Overlapping windows from an episode are correlated, and the number of valid starts decreases with horizon. Treating all windows as independent replications would substantially overstate precision. The held-out episodes are also used for validation logging, so these offline measurements are not an untouched final test set.

There are further comparison limits. WAM training uses horizon-length windows, excluding the final $H-1$ starts of each episode; behavioral cloning uses one-step windows. Thus the models do not see identical sets of supervised policy starts. Changing exploration also changes trajectory distributions, termination times, and normalization scales. A comparison between exploration conditions is suggestive, but does not isolate the effect of action diversity while holding everything else fixed.

\section{Archived results}
\subsection{State prediction is accurate relative to persistence}
With state inputs, the WAM's 15-step state NMSE is $0.00712\pm0.00038$, compared with $0.58243\pm0.01478$ for copying the initial state. The ratio of aggregate means is about 82. This establishes that the model predicts motion substantially better than persistence on the recorded trajectories. It does not establish an advantage over system identification, a state-space MLP, or a model-based controller.

The result also includes two static goal coordinates, which are directly supplied to the encoder. Low error on the six-dimensional average can conceal larger errors on moving quantities. Reporting separate position and velocity errors, together with physical units, is necessary for interpreting control relevance. In Section~\ref{sec:audit}, we use those separate quantities for the fresh baseline study.

\subsection{Action coverage changes the prediction problem}
\begin{figure}[t]
\centering
\includegraphics[width=\linewidth]{\figpath/research-coverage.pdf}
\caption{Archived neural prediction errors, means and population standard deviations over three runs. Expert-only data gives lower long-horizon factual error but higher error under random counterfactual actions. The two panels use different state dimensions; exploration conditions also have different trajectory distributions and normalization statistics. The curves demonstrate an evaluation gap, not a controlled causal estimate of exploration's benefit.}
\label{fig:coverage}
\end{figure}

The expert-only condition has a 15-step factual NMSE of $0.00469\pm0.00249$, lower on average than the standard exploratory condition. Its 10-step counterfactual NMSE is $0.11209\pm0.00499$, compared with $0.00718\pm0.00053$ at burst probability 0.3 (Figure~\ref{fig:coverage}). The latter ratio is approximately 15.6. A probability of 0.6 gives $0.00633\pm0.00229$ counterfactual error, with no corresponding improvement in factual error.

These results are consistent with a neural predictor relying too heavily on regularities of the behavior policy. They do not directly show that it ignores its action input: that claim would require action swaps, input ablation, or measured action sensitivity at fixed states. Nor do the results imply that expert data necessarily lacks enough information to identify the transition law. The affine experiment below provides a counterexample within this environment.

\subsection{Longer training horizons change long-range error}
\begin{table}[t]
\centering\small
\begin{tabular}{lrrr}
\toprule
Configuration & Factual $k=1$ & Factual $k=15$ & Counterfactual $k=10$\\
\midrule
Standard & .00188 & .00712 & .00718\\
One-step training & .00185 & .01100 & .01389\\
Ten-step training & .00295 & .00459 & .00569\\
No latent loss & .00164 & .00958 & .00698\\
No image loss & .00203 & .00743 & .00779\\
\bottomrule
\end{tabular}
\caption{Archived mean state NMSE across three runs. Factual errors average six dimensions; counterfactual errors average four. The table describes this configuration sweep, without significance or equivalence claims.}
\label{tab:ablation}
\end{table}

Ten-step training reduces the observed 15-step error while increasing one-step error (Table~\ref{tab:ablation}). Removing latent consistency increases the mean long-horizon factual error, but the counterfactual mean changes little. Removing image loss gives a similar mean to the standard condition, with a wider across-run spread at horizon 15: $0.00211$ versus $0.00038$. Three runs do not establish equivalence, and these measurements do not prove that either loss is unnecessary in a visual task. Here the encoder already receives the physical state.

\subsection{Decoded images reveal a different failure}
\begin{figure}[t]
\centering
\includegraphics[width=\linewidth]{\figpath/research-observations.pdf}
\caption{An archived held-out episode: observed frames, reconstructions from each current observation, and open-loop predictions conditioned on recorded actions. Rows use the same saved timestamps. The predicted row starts from one encoded observation; it does not receive the later frames. The imperfect rendering is retained as evidence rather than replaced with an idealized scene. This single episode is illustrative, not an aggregate image-quality result.}
\label{fig:images}
\end{figure}

The image decoder performs worse than persistence at every reported evaluation horizon. At one step its pixel MSE is $0.00642$, versus $0.00105$ for copying the current image. At 15 steps the values are $0.00730$ and $0.00500$. Figure~\ref{fig:images} shows the mismatch between recognizable scene structure and incomplete reconstructions.

A largely static background makes copying a strong pixel baseline. Small moving objects occupy relatively few pixels, so whole-image MSE also underweights their importance. These facts explain why image error and control relevance may disagree, but they do not excuse a failed comparison. The current decoder cannot support a claim of accurate visual imagination. Object-position error, goal retention, and foreground-weighted error would provide complementary diagnostics. Whether the decoder loses the distractor because it is irrelevant to control remains a hypothesis; its cause has not been isolated.

\subsection{Policy improvement depends on the input regime}
\begin{table}[t]
\centering\small
\begin{tabular}{lrr}
\toprule
Condition & WAM success (\%) & BC success (\%)\\
\midrule
400 episodes, state input & $100.0\pm0.0$ & $100.0\pm0.0$\\
25 episodes, 15 epochs & $87.0\pm7.8$ & $99.7\pm0.5$\\
25 episodes, 225 epochs & $96.3\pm1.2$ & $98.7\pm1.9$\\
400 episodes, pixel input & $80.3\pm12.0$ & $36.0\pm19.8$\\
\bottomrule
\end{tabular}
\caption{Archived closed-loop success: mean and population standard deviation across three runs, each evaluated on 100 episodes. The longer 25-episode run approximately matches the standard update budget; it does not match total computation between architectures. BC denotes the repository's VLA comparator.}
\label{tab:policy}
\end{table}

With state inputs, both policies succeed on all 300 evaluation episodes in the standard condition. This is a ceiling effect, not evidence that their behavior is identical. With 25 collection episodes and 15 epochs, auxiliary prediction coincides with worse control. Training for 225 epochs narrows the gap substantially (Table~\ref{tab:policy}), making optimization budget a plausible contributor to the short-run deficit. It does not establish the remaining gap's cause.

The pixel-input result is more interesting. The WAM outperforms BC in each of the three runs, by 44.3 percentage points on average, but the paired differences range from 5 to 74 points. This is a promising observation with substantial variability, not a stable estimate of a general world-model benefit.

Most importantly, WAM has supervised state targets that BC lacks. A policy trained with current-state reconstruction alone could benefit from the same localization signal without learning transitions. Future-state supervision, image reconstruction, latent consistency, and changes in the usable training windows are also entangled. The fresh pilot in Section~\ref{sec:pilot} begins to separate these effects. The archived comparison alone cannot establish that learned dynamics is the cause.

\section{Fresh simulator and system-identification audit}
\label{sec:audit}
To test whether the neural results reflect intrinsic task difficulty, we implemented a separate CPU experiment with five independent runs. It fits an affine transition predictor from state and action and compares it with an affine predictor that does not receive action. Each collection condition uses 60 training episodes and a separate bank of 20 held-out episodes, with a 60-step cap. Counterfactual tests branch from up to 512 held-out states per run under ten independent uniform actions. Both predictors use privileged state. These are diagnostic baselines for the transition problem, not alternatives to a visual encoder or a complete policy.

\subsection{A structured baseline recovers the dynamics}
The action-conditioned predictor fits a state increment from training-standardized features $[s_t,a_t,1]$ and adds it to $s_t$. Its position outputs are clipped to the arena after each recursive step, matching the known boundary rule. The action-free comparator uses $[s_t,1]$. Fitting uses training trajectories only, while factual and randomized-action evaluation use separate states. All collection conditions share the evaluation bank within a run. Errors are reported separately for position and normalized velocity in raw coordinates, avoiding condition-specific NMSE denominators.

Even deterministic expert data yields a full-rank action-conditioned design matrix in all five runs. The matrix has nine columns and rank nine, with condition numbers between approximately 13.5 and 16.2. The fitted action-conditioned model's 10-step counterfactual position MSE is about $4.3\times10^{-16}$, near the numerical floor of this test. The action-free model has much larger counterfactual position error, about $9.9\times10^{-3}$ at ten steps. Its one-step factual position MSE is about $2.1\times10^{-6}$.

This result changes the interpretation of the archived experiment. The neural model may benefit from exploratory trajectories, but the deterministic expert trajectories already identify the transition within this restricted model class. The clipped expert is piecewise affine, and variation across its operating regions can break the linear dependence that would arise from a single unconstrained linear feedback rule. The observed full rank is empirical evidence for these sampled datasets, not a theorem about every expert policy.

The affine model is favored by correct structural assumptions, full-state access, and the known clipping operation. Its numerical error must not be ranked directly against archived neural NMSE: the datasets, scales, and protocols differ. Its role is to establish that a matched structured baseline is necessary before calling this environment a demanding world-model benchmark.

\subsection{Observation ambiguity persists under imagination}
A paired simulator check holds agent position and both goals fixed, then assigns opposite velocities. The images are bit-identical, while applying the same action produces different next positions. A second check compares the environment's step function with its counterfactual simulator across 50 starts and 20 steps; they agree exactly after aligning the float32 starting state. These tests validate the intended observation ambiguity and transition oracle.

For an interior trajectory, two possible initial normalized velocities under identical subsequent controls satisfy
\begin{equation}
 \Delta u_{t+k}=0.8^k\Delta u_t.
\end{equation}
Thus uncertainty about initial velocity decays without any new observation. After five steps its squared contribution is scaled by $0.8^{10}\approx0.107$. A better future-velocity readout can therefore result from damping and known action inputs; it need not mean the model inferred the hidden initial velocity. In the archived offline test, those inputs are the recorded future actions, which also carry information about the expert's behavior.

This matters for the earlier pixel-input interpretation. The deployed policy re-encodes one frame each step and does not use these imagined latents. Improved future readout accuracy cannot by itself explain improved deployed control. Establishing that mechanism requires history-based controls, auxiliary-supervision controls, and evaluation on deliberately ambiguous scenes.

\section{A matched supervision pilot on Explorer}
\label{sec:pilot}
We next trained three pixel-input conditions on Northeastern's Explorer cluster: behavioral cloning alone, cloning with current-state prediction, and cloning with current and five-step future-state prediction. Each of five independent runs collects 200 episodes at burst probability 0.3, holds out 10\% by episode, and trains for 15 epochs. All conditions use identical five-step training starts, expert masks, encoder and policy initialization, minibatch order, and update counts within a run. Current and future state losses each have weight 0.5. Pixel reconstruction and latent consistency are excluded, so this is a targeted supervision experiment rather than a replication of the full archived WAM.

\begin{table}[t]
\centering\small
\begin{tabular}{lrr}
\toprule
Supervision & Mean $\pm$ SD & Range\\
\midrule
Actions only & $17.0\pm10.1$ & 0--27\\
+ Current state & $28.4\pm5.5$ & 24--38\\
+ Future state & $33.6\pm2.5$ & 30--36\\
\bottomrule
\end{tabular}
\caption{Fresh Explorer pilot: success percentages on 100 paired evaluation episodes per training run. The last row includes current-state supervision. SD is the sample standard deviation over five runs, unlike the population SD in the archived tables.}
\label{tab:pilot}
\end{table}

Current-state supervision improves the mean by 11.4 percentage points over cloning. Adding future-state prediction gives a further 5.2 points on average; the paired difference is positive in four of five runs and ranges from $-2$ to 9 points (Table~\ref{tab:pilot}). Thus current-state supervision accounts for a substantial part of the gain in this pilot, while future prediction adds a smaller, non-uniform benefit. The result is more informative than comparing the full auxiliary objective only against cloning.

Absolute success remains low, and this pilot uses half the archived collection budget. It neither establishes convergence nor explains the magnitude of the archived 44.3-point gap. It also lacks a history baseline and a state-free training objective. These measurements justify a larger, preregistered mechanism study; they do not yet establish a general advantage for learned dynamics. Per-episode outcomes, dataset hashes, source hashes, and configuration manifests are retained for all 15 trained models. The models act reactively throughout evaluation.

\section{What would establish a stronger research claim?}
The most useful next question is narrower than whether world models are helpful: under what observation and data-coverage conditions does auxiliary transition prediction improve control beyond equally supervised state estimation? A second question is whether decision-relevant counterfactual tests predict planning performance better than average state error. The pilot provides preliminary evidence for the first question; the second still requires a planning experiment.

A controlled study should first equalize policy training starts, expert-action labels, update counts, and evaluation episodes. It should compare BC with current-state supervision, one-step future-state supervision, and multistep prediction, both with and without privileged training targets. A short-history policy is essential because velocity ambiguity is introduced by the observation design. Parameter and wall-clock budgets should be reported alongside environment transitions.

Action coverage should then be varied with a fixed evaluation bank and a common metric scale. Holding out candidate actions at matched states separates action generalization from a change in the state distribution. Action-shuffled training and action-free predictors test whether transitions actually depend on the supplied controls. Affine identification, a nonlinear state-space model, and oracle dynamics bound the interpretation of latent-model results.

Finally, the model must influence decisions if the claim concerns planning. A common receding-horizon planner should use learned and oracle models under the same candidate budget and cost. Report control success, cost regret relative to the oracle candidate, model-call count, and latency. A predictor that reduces average error but ranks the useful candidates incorrectly is not an improved controller.

The point-mass task should remain a diagnostic setting. A publication-level robotics claim needs contact-rich or partially observed tasks with a reason for prediction to matter, multiple dynamics regimes, and comparisons to contemporary methods. A separate research plan accompanying the source specifies these experiments and their decision gates. They are proposed work, not completed results in this paper.

\section{Limitations and conclusion}
The study is limited by a single low-dimensional simulator, three archived neural runs, five fresh pilot runs, and incomplete original training provenance. The state-input task leaks goal identity through the selected target coordinates. The pixel-input task uses privileged supervision and lacks a memory baseline. The linear readout is jointly supervised, the image decoder underperforms persistence, and the forward model is not used for deployed planning. These are restrictions on what the results establish, not details that a change in writing style can remove.

Within that scope, mini-WAM is useful as an experimental instrument. It exposes a gap between factual and counterfactual prediction, a separation between state and image quality, and a possible benefit from auxiliary supervision under restricted observations. The fresh affine audit also supplies an instructive negative result: the neural counterfactual failure is not inevitable for the available expert data. A research contribution should explain these differences through controlled experiments. The current evidence supports that investigation, but not a claim that a new world-action method has already improved robotic decision-making.

\paragraph{Reproducibility and provenance.}
Archived measurements and figure arrays are retained under \texttt{docs/results/}; the new simulator audit records its full configuration, software versions, and source hashes separately. Environment images are actual renders, and neural reconstructions are preserved from the archived arrays. The accompanying research plan distinguishes completed diagnostics from planned training and robotics experiments. The source repository is \url{https://github.com/keivalya/mini-wam}.

\paragraph{Assistance disclosure.}
AI assistance was used to audit code, develop diagnostic experiments, check literature and venue information, and revise this manuscript. The results presented as fresh measurements were executed; archived results are labeled as such. Author review of the scientific claims and submission-specific disclosure remains necessary.

\begin{thebibliography}{9}
\bibitem[Alain and Bengio(2016)]{alain2016} G. Alain and Y. Bengio. Understanding intermediate layers using linear classifier probes. arXiv:1610.01644, 2016.
\bibitem[Chen et~al.(2026)]{chen2026} J. Chen et al. MiniWAM: Learning compact future targets for efficient world-action modeling. arXiv:2610.12194, 2026. \url{https://arxiv.org/abs/2610.12194}.
\bibitem[de Haan et~al.(2019)]{dehaan2019} P. de Haan, D. Jayaraman, and S. Levine. Causal confusion in imitation learning. NeurIPS, 2019. \url{https://arxiv.org/abs/1905.11979}.
\bibitem[Ha and Schmidhuber(2018)]{ha2018} D. Ha and J. Schmidhuber. Recurrent world models facilitate policy evolution. NeurIPS, 2018. \url{https://arxiv.org/abs/1809.01999}.
\bibitem[Hafner et~al.(2020)]{hafner2020} D. Hafner, T. Lillicrap, J. Ba, and M. Norouzi. Dream to control: Learning behaviors by latent imagination. ICLR, 2020. \url{https://arxiv.org/abs/1912.01603}.
\bibitem[Hansen et~al.(2024)]{hansen2024} N. Hansen, H. Su, and X. Wang. TD-MPC2: Scalable, robust world models for continuous control. ICLR, 2024. \url{https://arxiv.org/abs/2310.16828}.
\bibitem[Qiu et~al.(2026)]{qiu2026} J. Qiu, Z. Chen, H. Cao, J. Shi, J. Huo, and Y. Gao. AD-WM: Action-discriminative world models for counterfactual model predictive control. arXiv:2609.30264, 2026. \url{https://arxiv.org/abs/2609.30264}.
\bibitem[Zhu et~al.(2025)]{zhu2025} C. Zhu et al. Unified world models: Coupling video and action diffusion for pretraining on large robotic datasets. Robotics: Science and Systems, 2025. \url{https://doi.org/10.15607/RSS.2025.XXI.015}.
\end{thebibliography}
\end{document}
